\documentclass[11pt,a4paper]{article}
\usepackage[italian]{babel}
\usepackage[margin=2.5cm]{geometry}
\usepackage{parskip}
\usepackage{amsmath, amssymb, amsthm}
\usepackage{booktabs}
\usepackage{array}
\usepackage{xcolor}
\usepackage{needspace}
\usepackage{enumitem}
\usepackage{listings}
\usepackage{tikz}
\usepackage[most]{tcolorbox}
\usepackage[hidelinks]{hyperref}
\usetikzlibrary{decorations.pathreplacing, shapes.geometric, arrows.meta, positioning}

% Niente vedove/orfane: i paragrafi non lasciano righe isolate a fine/inizio pagina
\widowpenalty=10000
\clubpenalty=10000

% Elenchi più compatti
\setlist{itemsep=2pt, parsep=0pt, topsep=4pt, partopsep=0pt}

% --- Riquadri con bordo nero, senza numero (si spezzano tra due pagine) ---
% Uso: \begin{definizione} ... \end{definizione}
%      \begin{teorema}[Nome] ... \end{teorema}  (il nome è facoltativo)
\tcbset{nota/.style={enhanced jigsaw, breakable, colback=white,
  colframe=black, boxrule=0.5pt, sharp corners}}
\NewTColorBox{definizione}{}{nota,
  before upper={\textbf{Definizione.}\ }}
\NewTColorBox{teorema}{o}{nota,
  before upper={\textbf{Teorema\IfValueT{#1}{ (#1)}.}\ }}
\NewTColorBox{proposizione}{o}{nota,
  before upper={\textbf{Proposizione\IfValueT{#1}{ (#1)}.}\ }}
\NewTColorBox{proprieta}{o}{nota,
  before upper={\textbf{Proprietà\IfValueT{#1}{ (#1)}.}\ }}
\NewTColorBox{assioma}{o}{nota, boxrule=1pt,
  before upper={\textbf{Assioma\IfValueT{#1}{ (#1)}.}\ }}

% --- Ambienti semplici (senza riquadro) ---
% Il titolo sta in cima al blocco, anche se il contenuto inizia con un riquadro o
% con codice; e non resta mai da solo in fondo a una pagina.
% Uso: \begin{esempio}[Titolo facoltativo] ... \end{esempio}
\newcommand{\newrunin}[2]{%
  \NewDocumentEnvironment{#1}{o}{%
    \par\Needspace{4\baselineskip}\addvspace{\medskipamount}\noindent
    \textbf{#2}\IfValueT{##1}{ (##1)}\textbf{.}\ \ignorespaces}%
  {\par\addvspace{\medskipamount}}}
\newrunin{esempio}{Esempio}
\newrunin{esercizio}{Esercizio}
\newrunin{osservazione}{Osservazione}

% V e F nelle tavole di verità
\newcommand{\V}{\textsf{V}}
\newcommand{\F}{\textsf{F}}

% --- Codice C ---
% Uso: \begin{codice} ... \end{codice}      codice C numerato
%      \begin{comando} ... \end{comando}    comandi da terminale
%      \begin{risultato} ... \end{risultato} output di un programma
%      \lstinline|printf("ciao");|           codice dentro al testo
\lstdefinestyle{c}{
  language=C,
  basicstyle=\ttfamily\small,
  keywordstyle=\color{blue}\bfseries,
  commentstyle=\color{gray}\itshape,
  stringstyle=\color{red!70!black},
  numbers=left,
  numberstyle=\tiny\color{gray},
  numbersep=8pt,
  columns=flexible,
  keepspaces=true,
  breaklines=true,
  showstringspaces=false,
  tabsize=4,
  morekeywords={include, define},
  % lettere accentate nei commenti
  literate={à}{{\`a}}1 {è}{{\`e}}1 {é}{{\'e}}1 {ì}{{\`i}}1
           {ò}{{\`o}}1 {ù}{{\`u}}1 {È}{{\`E}}1 {À}{{\`A}}1
}
\lstset{style=c}
\newtcblisting{codice}{listing only, nota, left=6mm,
  listing options={style=c}}
\newtcblisting{comando}{listing only, nota,
  listing options={basicstyle=\ttfamily\small, numbers=none, language={},
    columns=flexible, keepspaces=true, breaklines=true}}
\newtcblisting{risultato}{listing only, enhanced jigsaw, breakable,
  colback=black!5, colframe=black, boxrule=0.5pt, sharp corners,
  listing options={basicstyle=\ttfamily\small, numbers=none, language={},
    columns=flexible, keepspaces=true, breaklines=true}}

% --- Diagrammi di flusso (simboli standard) ---
\tikzset{
  terminale/.style = {ellipse, draw, minimum width=2.5cm, minimum height=0.9cm},
  io/.style        = {trapezium, trapezium left angle=70, trapezium right angle=110,
                      draw, minimum width=2.5cm, minimum height=0.9cm},
  processo/.style  = {rectangle, draw, minimum width=2.5cm, minimum height=0.9cm},
  decisione/.style = {diamond, draw, aspect=2, minimum width=2.5cm},
  freccia/.style   = {thick, -{Stealth}}
}

\title{Inclusione-esclusione e conteggio}
\author{Marco Cerbella}
\date{Lezione 8 -- 13 ottobre 2026}

\begin{document}
\maketitle

\section{Principio di inclusione-esclusione}

\begin{proposizione}
Siano $A, B$ due insiemi finiti. Allora
\[
|A \cup B| = |A| + |B| - |A \cap B|.
\]
\end{proposizione}

\emph{Dimostrazione.} Sommando $|A|$ e $|B|$ gli elementi di $A \cap B$ vengono contati due volte. Perciò $|A \cup B| + |A \cap B| = |A| + |B|$, da cui la formula. \hfill $\square$

\begin{esempio}
Quanti numeri naturali positivi $\le 43$ sono divisibili per $4$ oppure per $5$?
Sia $A = \{n \in \mathbb{N} : n \le 43,\ 4 \mid n\}$ e $B = \{n \in \mathbb{N} : n \le 43,\ 5 \mid n\}$. Allora $|A| = 10$ e $|B| = 8$. L'intersezione è
\[
A \cap B = \{n \le 43 : 5 \mid n \text{ e } 4 \mid n\} = \{n \le 43 : 20 \mid n\} = \{20, 40\},
\]
quindi $|A \cap B| = 2$ e $|A \cup B| = 10 + 8 - 2 = 16$.
\end{esempio}

\begin{proposizione}
Siano $A, B, C$ tre insiemi finiti. Allora
\[
|A \cup B \cup C| = |A| + |B| + |C| - |A \cap B| - |A \cap C| - |B \cap C| + |A \cap B \cap C|.
\]
\end{proposizione}

\begin{esempio}[Evento sportivo]
Un evento ha tre gare e $20$ atleti partecipanti. Sia $A$ l'insieme degli atleti della gara $1$, $B$ della gara $2$, $C$ della gara $3$. Si sa che $|A| = 6$, $|B| = 10$, $|C| = 13$, $|A \cap B| = 4$, $|A \cap C| = 3$, $|B \cap C| = 5$. Quanti atleti partecipano a tutte e tre le gare?
Il numero totale di atleti è $|A \cup B \cup C| = 20$. Ricavando $|A \cap B \cap C|$ dalla formula:
\[
|A \cap B \cap C| = |A \cup B \cup C| - |A| - |B| - |C| + |A \cap B| + |A \cap C| + |B \cap C| = 20 - 6 - 10 - 13 + 4 + 3 + 5 = 3.
\]
\end{esempio}

\section{Contare le funzioni tra insiemi finiti}

Siano $A, B$ insiemi finiti con $|A| = n$ e $|B| = k$.

\begin{itemize}
    \item \emph{Funzioni da $A$ a $B$}: per ogni elemento di $A$ si sceglie un'immagine tra i $k$ elementi di $B$. Sono le disposizioni con ripetizioni: $D'_{k,n} = k^n$.
    \item \emph{Funzioni biunivoche} (possibili se $n = k$): sono le permutazioni di $n$ oggetti, $n! = n \cdot (n-1) \cdots 2 \cdot 1$.
    \item \emph{Funzioni iniettive} (possibili se $n \le k$): sono le disposizioni semplici, $D_{k,n} = k(k-1)\cdots(k-n+1)$.
    \item \emph{Funzioni suriettive} (possibili se $n \ge k$):
    \begin{itemize}
        \item se $|B| = 1$ c'è la sola funzione costante, che è suriettiva: $1$ funzione;
        \item se $|B| = 2$ le funzioni sono $2^n$; non sono suriettive le $2$ funzioni costanti, quindi le suriettive sono $2^n - 2$;
        \item se $|B| = 3$ le funzioni sono $3^n$. Quelle che lasciano fuori almeno un elemento di $B$ sono $3 \cdot 2^n - 3$: il termine $3 \cdot 2^n$ conta, per ciascuno dei $3$ elementi, le $2^n$ funzioni che lo escludono, ma le $3$ funzioni costanti (che escludono due elementi) sono state contate due volte, quindi se ne sottraggono $3$. Le suriettive sono
        \[
        3^n - (3 \cdot 2^n - 3) = 3^n - 3 \cdot 2^n + 3.
        \]
    \end{itemize}
\end{itemize}

\section{Cardinalità dell'insieme delle parti}

\begin{proposizione}
Se $A$ è un insieme finito con $|A| = n$, l'insieme delle parti $\mathcal{P}(A)$ ha cardinalità $2^n$.
\end{proposizione}

\emph{Dimostrazione} per induzione su $n$. Sia $P(n)$: se $|A| = n$ allora $|\mathcal{P}(A)| = 2^n$.

\emph{Caso base} $n = 0$: $A = \emptyset$ e $\mathcal{P}(\emptyset) = \{\emptyset\}$, che ha $1 = 2^0$ elementi.

\emph{Passo induttivo.} Ipotesi: $P(k)$, cioè un insieme con $k$ elementi ha $2^k$ sottoinsiemi. Sia $|A| = k + 1$; scegliamo $\bar a \in A$. I sottoinsiemi $X \subseteq A$ sono di due tipi:
\begin{itemize}
    \item $\bar a \notin X$: allora $X \subseteq A \setminus \{\bar a\}$, che ha $k$ elementi, e per ipotesi induttiva questi $X$ sono $2^k$;
    \item $\bar a \in X$: allora $X = \{\bar a\} \cup Y$ con $Y \subseteq A \setminus \{\bar a\}$; di nuovo i possibili $Y$ sono $2^k$.
\end{itemize}
Quindi $|\mathcal{P}(A)| = 2^k + 2^k = 2 \cdot 2^k = 2^{k+1}$. \hfill $\square$

\section{Esercizi di conteggio}

\begin{esempio}[Carte]
Un mazzo ha $40$ carte ($4$ semi, $10$ valori per ogni seme).
\begin{enumerate}
    \item Quante mani di $5$ carte? Combinazioni semplici: $\binom{40}{5} = \dfrac{40!}{5!\,35!}$.
    \item Quante mani con $4$ assi? Le $4$ carte sono fissate, la quinta è una delle altre $36$: $\binom{36}{1} = 36$.
    \item Quante mani con $4$ carte dello stesso valore? Si sceglie il valore ($10$ modi) e poi l'ultima carta ($36$ modi): $10 \cdot 36$.
    \item Quante mani con \emph{esattamente} $2$ assi? Si scelgono i $2$ assi e poi $3$ carte che non sono assi: $\binom{4}{2}\binom{36}{3}$.
    \item Quante mani con \emph{almeno} $2$ assi? Si sommano i casi con esattamente $2$, $3$ o $4$ assi:
    \[
    \binom{4}{2}\binom{36}{3} + \binom{4}{3}\binom{36}{2} + \binom{4}{4}\binom{36}{1}.
    \]
    \item Quante mani con $2$ coppie di valori diversi (più una quinta carta)? Si scelgono i due valori ($\binom{10}{2}$), poi $2$ carte per il primo valore ($\binom{4}{2}$), $2$ per il secondo ($\binom{4}{2}$) e la quinta carta tra le $32$ rimaste con altri valori:
    \[
    \binom{10}{2}\binom{4}{2}\binom{4}{2}\binom{32}{1}.
    \]
\end{enumerate}
\end{esempio}

\begin{esempio}[Panchina]
In quanti modi $3$ maschi e $3$ femmine possono sedersi su una panchina?
\begin{enumerate}
    \item Senza vincoli: permutazioni di $6$ oggetti, $6!$.
    \item Se i maschi stanno tutti vicini e le femmine tutte vicine: $2$ ordini dei due gruppi, $3!$ per i maschi, $3!$ per le femmine: $2 \cdot 3! \cdot 3!$.
    \item Se solo i maschi stanno vicini: i $3$ maschi formano un blocco che, insieme alle $3$ femmine, dà $4$ oggetti da permutare ($4!$), poi si ordinano i maschi nel blocco ($3!$): $4! \cdot 3!$.
    \item Se maschi e femmine si alternano: si sceglie chi inizia ($2$), poi si ordinano i maschi ($3!$) e le femmine ($3!$): $2 \cdot 3! \cdot 3!$.
\end{enumerate}
\end{esempio}

\begin{esempio}[Urna]
Un'urna contiene $20$ palline bianche e $10$ nere. Si fanno $5$ estrazioni e si vuole ottenere $3$ bianche e $2$ nere.
\begin{enumerate}
    \item \emph{Senza reimmissione}: combinazioni semplici, $\binom{20}{3}\binom{10}{2}$.
    \item \emph{Con reimmissione} (si contano gli insiemi di palline scelte): combinazioni con ripetizioni,
    \[
    C'_{20,3}\cdot C'_{10,2} = \binom{22}{3}\binom{11}{2},
    \]
    perché $C'_{n,k} = C_{n+k-1,k}$.
    \item \emph{Senza reimmissione, palline indistinguibili e si tiene conto dell'ordine} (sequenze come $BBBNN$, $NBBNB, \dots$): sono permutazioni con ripetizioni, $\dfrac{5!}{3!\,2!}$.
\end{enumerate}
\end{esempio}

\section{Un esercizio di induzione}

\begin{esercizio}
Dimostrare per induzione che $25 \mid 2^{5n} - 7^n$ per ogni $n \ge 1$.
\end{esercizio}

\emph{Soluzione.} Sia $P(n)$: $25 \mid 2^{5n} - 7^n$.

\emph{Caso base} $n = 1$: $2^5 - 7 = 32 - 7 = 25$ e $25 \mid 25$.

\emph{Passo induttivo.} Ipotesi: $25 \mid 2^{5k} - 7^k$, cioè esiste $h \in \mathbb{Z}$ con $2^{5k} - 7^k = 25h$, ossia $2^{5k} = 25h + 7^k$. Dobbiamo dimostrare $25 \mid 2^{5(k+1)} - 7^{k+1}$. Si ha
\begin{align*}
2^{5k+5} - 7^{k+1} &= 2^5 \cdot 2^{5k} - 7 \cdot 7^k = 2^5 (25h + 7^k) - 7 \cdot 7^k \\
&= 25 \cdot 2^5 h + 7^k (2^5 - 7) = 25 \cdot 2^5 h + 7^k \cdot 25 \\
&= 25\,(2^5 h + 7^k),
\end{align*}
quindi $25 \mid 2^{5(k+1)} - 7^{k+1}$, che era la tesi. \hfill $\square$

\end{document}
